% \documentclass[10pt,preprint]{sigplanconf}
\documentclass[10pt]{article}

\begin{document}

% \special{papersize=8.5in,11in}
% \setlength{\pdfpageheight}{\paperheight}
% \setlength{\pdfpagewidth}{\paperwidth}
% \setlength{\hoffset}{-1in}
% \setlength{\voffset}{-1in}
\setlength{\topmargin}{1in}
\setlength{\headheight}{0pt}
\setlength{\headsep}{0pt}
\setlength{\evensidemargin}{.75in}
\setlength{\oddsidemargin}{.75in}
\newdimen{\standardtextwidth}
\setlength{\standardtextwidth}{42pc}
\setlength{\textwidth}{\standardtextwidth}
\setlength{\topskip}{8pt}
\setlength{\columnsep}{2pc}
\setlength{\textheight}{54.5pc}
\setlength{\parskip}{0pt}
\setlength{\parindent}{12pt}
\setlength{\lineskip}{.5pt}
\setlength{\lineskiplimit}{\lineskip}
\frenchspacing
\pretolerance = 400
\tolerance = \pretolerance
\setlength{\emergencystretch}{5pt}
\clubpenalty = 100000
\widowpenalty = 100000
\setlength{\hfuzz}{.5pt}
\setlength{\marginparwidth}{36pt}
\setlength{\marginparsep}{2pt}
\setlength{\marginparpush}{8pt}
\setlength{\skip\footins}{8pt plus 3pt minus 1pt}
\setlength{\footnotesep}{9pt}

\date{}

% Communications Design Group, SAP Labs, San Francisco, CA, USA}
% \maketitle
\sloppy
\begin{center}
{\large Cover Letter for Revisions to Paper}\\
{\Large Implementing Record and Refinement for\\
  Debugging Timing-dependent Communication}  
\end{center}

\vskip2\baselineskip

\section*{Response to Minor Revision Comments}
We again thank the reviewers for taking the time and further commenting on our
work. We have addressed the minor remarks regarding spelling and grammar. We
detail how we addressed the two larger remarks below.

\paragraph{Phrasing of Contributions in Abstract and the list of Contributions}
Reviewer 3 remarked that the contributions in the abstract are not phrased as
such. We have fixed this, and now more clearly state what we show. Furthermore,
we have extended the description of the last contribution to clearly show that
Peek-at-Talk is a debugger integrated into our Path Tools framework.

\paragraph{Evaluation}
Reviewer 3 also asked about a comment regarding our user study that went
unaddressed in the major revision. We apologize for this mistake. We have now
added details about our preliminary findings in the user study. While we did not
have the resources to conduct a larger study, our preliminary findings indicate
the tool is useful for actual developers in some cases. We have added a few
sentences describing these findings.

% \section*{Response to Major Revision Comments}
% Thanks to the reviewers for their detailed and useful comments. This letter
% summarizes our changes in response to the more general comments. We have also
% addressed various smaller and more direct comments that dealt with sentences
% that were hard to understand, misspelled words, as well as grammatical mistakes.
% % Besides listing the raised issues and showing our changes, we have
% % also attached a full-text diff at the end of this document.

% \paragraph{Approach and Implementation Details}
% A major concern in particular of reviewer 1 was that we present our
% implementation in much detail, but do not present the core ideas of our approach
% separately. Since our approach has grown out of a concrete problem and an
% implementation that we developed to address this problem, we have only mildly
% edited the implementation details, but we have added a section that distills our
% approach and concisely presents the core ideas that we think are relevant for an
% attempt to implement our approach in a different system.

% We have also re-organized our list of contributions to make it clearer what the
% contributions with respect to the approach are, and what contributions are with
% respect to the implementation.

% \paragraph{Clarity about Constraints of our Solution and Application to Example}
% The reviewers remarked that we should be clearer already in the introduction
% about the kinds of scenarios that we present our work for. We have added a
% sentence in the introduction to clarify that it is useful for cases where
% variances in network events cause failures. We hope this also clarifies the
% constraints on scenarios in which we think our approach is useful. We have also
% discussed this in the newly added Section 5.

% \paragraph{Interactions with our tool}
% The reviewers asked that we summarize the interactions the developer has with
% our tool to address the bug we present in our example. Fixing the bug in the
% running example is spread over the entirety of section 4, and we agree that it
% was thus hard to find. We have added a paragraph at the introduction of section
% 4 to summarize the interactions that we show in the rest of the section.

% \paragraph{Using COP for wrapping the core library}
% The reviewers remarked that using COP to wrap the core library is not a strong
% contribution in itself. While we still think it is worthwhile to describe it so
% another researcher attempting to implement a similar system has an idea about
% how to implement the recording, traffic shaping, and refinement of data, we have
% removed it as a contribution.

% We have clarified that COP is available for many programming systems, and on
% many of these systems it can wrap around all methods required for our approach.

% \paragraph{Details about the Path Tools framework}
% The reviewers raised some questions about how much detail is required about the
% Path Tools framework. Since Peek-At-Talk integrates deeply with the Path Tools,
% and the idea of guiding the developer through multiple debugging steps at
% different levels of abstraction comes from our Path Tools framework, we think an
% overview of that framework is useful for future implementers of similar
% tools. However, we have removed some details about other Path Tools extensions
% such as ``dynamic contracts'' and ``inductive analysis''.

% % \paragraph{SHARE VM}
% % SHARE VM
% % 	- I could not figure out how the PathMap tool is structured. I have some green and white boxes. Green boxes are tests that have passed. White boxes represent uncovered methods. I could not figure out what the size and their position depends on. The questions I immediately bump into are: What makes a method square or rectangular. I would expect a tool tip info that indicate this to me. If the height and the width correspond to some kind of metrics (as promoted in Polymetric Views, by Lanza et al.), and I need to have at least some minimal indication what they correspond to.
% % 	- I was able to modify a test to make it fail. I wanted to see how Path Finder behaves in that case.
% % 	- I have raised an exception "subscript is out of bounds: 11", from THMonoTypeBucket>>memberTypeFor:in: by opening a PathView
% % 	- I then wanted to try our PathMap on my own unit tests. I therefore created a new test (I am confortable with Smalltalk programming). But I could not see how to lunch PathMap on my own test. I opened the PathBrowser, hoping it will give me an option to lunch PathMap. But when I selected my unit test, I got an error: MessageNotUnderstood: ClassOrganizer>>invariant. The top of the stack was SBrowser>>invariantsButtonColor . Pity, I wanted to try how the back-in-time debugger on my own test scenario :-( 
% % 	- The screencast was of a little use unfortunately. Looking at it on a remote image is not great due to the lag. I have hope to see it on a youtube or vimeo channel, but I have been unlucky.

% \paragraph{Related work for back-in-time debugging and record and replay}
% We have extended the related work to contrast our work with record and replay as
% well as back in time debuggers. The main advantage is the limited overhead
% required for our recording, because rather than recording all data at once, we
% incrementally record slices of data over multiple runs.

% \paragraph{Issues with Diagnosis Scripts}
% The reviewers asked about possible issues with diagnosis scripts, forcing
% different schedules pro-actively, and if the scripts themselves can lead to
% denial of service. We have extended our discussion of these scripts (also in the
% new section 5). The scripts must of course be run in a way to not cause a denial
% of service, and especially for more complex scenarios, it may take a very long
% time to collect enough data for our analysis to work. However, due to the low
% overhead of our approach, we argue that the script can easily run for days
% before it has collected enough data and proceeding with the analysis.

% We have also clarified that even in the open world case, our script can at least
% find anomalies caused by timing dependencies between the client and the server,
% even though code on non-cooperating clients cannot be further inspected.


\end{document}

%%% Local Variables:
%%% mode: latex
%%% TeX-master: t
%%% End:
